\begin{table}[H]
\centering
\caption{Nama Daerah Perbandingan Klaster Biplot vs Kategori IPM (UNDP) (2022)} \label{tab:app1crossname2022}
\begin{adjustbox}{max width=\textwidth}
\begin{tabular}{|>{\centering\arraybackslash}p{4cm}|p{10cm}|}
\hline
\multicolumn{1}{|c|}{\textbf{Kategori}} & \multicolumn{1}{c|}{\textbf{Daftar Wilayah}} \\
\hline
Biplot 1 (Rendah) \& IPM UNDP (Tinggi) & Kep. Seribu, Tangerang, Kota Cilegon, Kota Serang \\
\hline
Biplot 1 (Rendah) \& IPM UNDP (Sedang) & Pandeglang, Lebak, Serang \\
\hline
Biplot 2 (Sedang) \& IPM UNDP (Tinggi) & Bogor, Bandung, Ciamis, Kuningan, Cirebon, Sumedang, Purwakarta, Karawang, Kota Cirebon, Kota Tasikmalaya, Kota Banjar \\
\hline
Biplot 2 (Sedang) \& IPM UNDP (Sedang) & Sukabumi, Cianjur, Garut, Tasikmalaya, Majalengka, Indramayu, Subang, Bandung Barat, Pangandaran \\
\hline
Biplot 3 (Tinggi) \& IPM UNDP (Sangat Tinggi) & Kota Jakarta Selatan, Kota Jakarta Timur, Kota Jakarta Pusat, Kota Jakarta Barat, Kota Jakarta Utara, Kota Bandung, Kota Bekasi, Kota Depok, Kota Tangerang Selatan \\
\hline
Biplot 3 (Tinggi) \& IPM UNDP (Tinggi) & Bekasi, Kota Bogor, Kota Sukabumi, Kota Cimahi, Kota Tangerang \\
\hline
\end{tabular}
\end{adjustbox}
\end{table}